% Slide 2: What Is Spec-Driven Development?
% Visual: A pipeline diagram showing Idea -> Spec -> Code -> Ship, with the spec
%   stage expanded and emphasized. Below, a comparison: "Spec = the new compiler".
% Talking Points:
%   - SDD is writing a formal, structured specification BEFORE any code is written.
%   - The spec is not documentation -- it is the primary engineering artifact.
%   - Think of writing the spec as "compiling your intent" into machine-readable form.
% Presenter Script:
%   "So what is Spec-Driven Development? It's deceptively simple: you write a
%    rigorous, structured specification -- in Markdown, not UML -- before a single
%    line of code exists. This is NOT a requirements doc that sits in Confluence
%    and rots. This IS the engineering artifact. It defines functional requirements,
%    API contracts, data models, edge cases, and acceptance criteria. Here's the
%    mental model I want you to take away: writing the spec IS compiling. You are
%    compiling your intent from the messy human language in your head into a
%    precise, structured format that an AI agent can execute. If your spec doesn't
%    compile -- if it's ambiguous, incomplete, or contradictory -- the output code
%    will have bugs. Just like a real compiler."

\section{Spec-Driven Development}

\begin{frame}{What Is Spec-Driven Development?}
  \begin{center}
    \begin{tikzpicture}[scale=0.8, transform shape,
        stage/.style={rectangle, draw, rounded corners, minimum width=2.5cm,
                      minimum height=1.2cm, font=\scriptsize\bfseries, align=center},
        arrow/.style={->, >=stealth, thick}]

      \node[stage, fill=keylimegray!15] (idea) at (0,0) {\faIcon{lightbulb}\\Idea};
      \node[stage, fill=rhred!15, draw=rhred, line width=1.5pt] (spec) at (3.5,0) {\faIcon{file-alt}\\Spec};
      \node[stage, fill=keylimeblue!15] (code) at (7,0) {\faIcon{code}\\Code};
      \node[stage, fill=keylimegreen!15] (ship) at (10.5,0) {\faIcon{rocket}\\Ship};

      \draw[arrow, keylimegray] (idea) -- (spec);
      \draw[arrow, rhred] (spec) -- (code);
      \draw[arrow, keylimeblue] (code) -- (ship);

      % Expansion below spec
      \node[rectangle, draw, rounded corners, fill=rhred!5, text width=6cm,
            font=\tiny, inner sep=0.2cm, below=0.8cm of spec] (detail) {
        \textbf{The Spec contains:}\\
        Functional requirements, API contracts, data models,\\
        edge cases, error handling, acceptance criteria
      };
      \draw[dashed, rhred] (spec) -- (detail);
    \end{tikzpicture}
  \end{center}

  \vspace{0.3cm}

  \begin{columns}[T]
    \begin{column}{0.48\textwidth}
      \begin{block}{What It Is}
        \scriptsize
        \begin{itemize}
          \item A \textbf{structured Markdown document}
          \item The \textbf{primary engineering artifact}
          \item Machine-readable intent
          \item Versioned alongside code
        \end{itemize}
      \end{block}
    \end{column}
    \begin{column}{0.48\textwidth}
      \begin{alertblock}{What It Is NOT}
        \scriptsize
        \begin{itemize}
          \item A Confluence doc that rots
          \item A Jira ticket with two sentences
          \item UML diagrams nobody reads
          \item An afterthought or formality
        \end{itemize}
      \end{alertblock}
    \end{column}
  \end{columns}

  \vspace{0.3cm}

  \begin{center}
    \colorbox{rhred!10}{\parbox{0.8\textwidth}{\centering\scriptsize
      \textbf{Mental model:} Writing the spec = \textbf{compiling your intent}.\\
      Ambiguous spec = compiler errors. Clear spec = working code.
    }}
  \end{center}
\end{frame}
